\documentclass[11pt]{article}
\usepackage[margin=2.5cm]{geometry}
\usepackage{amsmath,amssymb}
\usepackage{booktabs}
\usepackage{array}
\usepackage{xcolor}
\usepackage{hyperref}
\usepackage{listings}
\usepackage{microtype}

\newcommand{\E}{\mathbb{E}}
\newcommand{\code}[1]{\texttt{#1}}
\let\origunderscore\_
\renewcommand{\_}{\origunderscore\allowbreak}
\newcommand{\deff}{d_{\mathrm{eff}}}

\lstset{basicstyle=\ttfamily\small, breaklines=true}
\emergencystretch=3em

\title{Effective dimensionality $\deff(x,\tau)$ from MGD's fitted $\theta_t$}
\author{Generated with Claude Code, from code-reading and debugging sessions,
  2026-09-23 to 2026-09-24}
\date{}

\begin{document}
\maketitle

\section*{Purpose}

Guth et al.\ define a local effective dimensionality $\deff(x,\tau)$ of a
distribution's support, at a point $x$ and noise scale $\tau$, from the
score of the variance-exploding (VE) noised distribution at $x$. This note
derives $\deff$ directly in terms of MGD's own fitted quantities --
$\theta_t$, $\varphi$, $\nabla\varphi$ -- so it can be computed as pure
post-processing of an already-completed MGD run, with no new
sampling/fitting. Implementation: \code{codes/effective\_dimension.py}.
Notebook: \code{turbulence/effective\_dimension\_deff.ipynb}.

\section{Guth's $\deff$: recap}

Let $q_\tau$ be the law of $y = x + \sqrt\tau\,\xi$, $\xi\sim\mathcal N(0,I_d)$,
$x$ drawn from the data. Two equivalent definitions:

\begin{align}
  \text{energy form:} \qquad
  \deff(x,\tau) &= d - \partial_{s}\, \E_y\!\big[U(y,\tau)\mid x\big], \qquad s = \tfrac12\log\tau,
  \label{eq:energy-def}\\[2mm]
  \text{denoiser form:} \qquad
  \deff(x,\tau) &= \frac1\tau\,\E_y\!\big[\lVert x - \E[x\mid y]\rVert^2 \mid x\big],
  \label{eq:denoiser-def}
\end{align}
where $U(\cdot,\tau) = -\log q_\tau(\cdot)$ and, by Tweedie's formula,
$\E[x\mid y] = y + \tau\,\nabla_y\log q_\tau(y)$. Both reduce to $d$ for an
isotropic Gaussian and to (less than) $d$ when the data lies near a
lower-dimensional manifold, at scales $\tau$ finer than the manifold's local
curvature.

\section{MGD's schedule and sign convention}

MGD's Cos stochastic interpolant (\code{sde\_routines.py}, \code{compute\_}
\code{interpolant}, \code{interpolant='Cos'} -- the only schedule
\code{forward\_regularised} supports) is
\begin{equation}
  I_t = \cos(\alpha_t)\,x_0 + \sin(\alpha_t)\,x_1, \qquad \alpha_t = \tfrac{\pi t}{2},
  \qquad x_0\sim\mathcal N(0,I_d),\ x_1 = x\ \text{(data)}.
\end{equation}

\paragraph{Sign convention.} This codebase fits $\theta_t$ for
$p_\theta(x)\propto\exp(+\theta^\top\varphi(x))$ -- the \emph{opposite} of
the MGD paper's $\exp(-\theta^\top\varphi(x))$. This is audited directly
against the code in \code{codes/utils\_entropy.py} (``Sign convention
audit''): the corrector step in \code{iteration\_step\_projection} is
$x_{k+1} = y_k + \nabla\varphi^\top\theta_{\mathrm{raw}}$ (a \emph{plus}),
giving a continuous-time drift $\sigma^2\nabla\varphi^\top\theta_t$ that
must equal $\sigma^2\nabla\log p_t$ for the transport to preserve $p_t$, so
\begin{equation}
  \nabla_x \log p_t(x) \;=\; +\,\theta_t^\top \nabla\varphi(x).
  \label{eq:score-sign}
\end{equation}
Independently corroborated by $dH_k = -\theta_k^\top \dot m_t$ (the entropy
increment already computed in \code{iteration\_step\_projection}) and by
$\log Z_{\theta_1} \gtrsim H_{\mathrm{bound}} + \theta_1^\top m_1$ (MGD
Eq.~8/30 as implemented in \code{utils\_entropy.log\_Z\_bound}). Every
formula below is re-derived directly in this $(+)$ convention -- not by
substituting $\theta_{\mathrm{paper}} = -\theta_{\mathrm{code}}$ into the
paper's stated formulas.

\section{The VE rescaling}

Define $y_t := I_t/\sin(\alpha_t) = x_1 + \cot(\alpha_t)\,x_0$. Conditional
on $x_1 = x$,
\begin{equation}
  y_t \mid x \;\sim\; \mathcal N\!\big(x,\ \tau(t)\,I_d\big), \qquad
  \tau(t) = \cot^2(\alpha_t),
  \label{eq:tau-of-t}
\end{equation}
monotonically decreasing from $\tau(0)=\infty$ to $\tau(1)=0$. So MGD's Cos
schedule already sweeps every VE noise scale, and $t_{\mathrm{reg}}$ (the
grid $\theta_t$ is fit on) IS a $\tau$-grid via Eq.~\ref{eq:tau-of-t}.

Write $q_\tau$ (Guth's VE marginal, density in $y$) in terms of $p_t$ (MGD's
marginal, density in $I_t$) by the change of variables $I_t = \sin(\alpha_t)\,y$:
\begin{equation}
  q_\tau(y) = \sin^d(\alpha_t)\; p_t\big(\sin(\alpha_t)\,y\big).
\end{equation}
Differentiating and using Eq.~\ref{eq:score-sign} (evaluated at $I_t = \sin(\alpha_t) y$):
\begin{equation}
  \nabla_y \log q_\tau(y) \;=\; \sin(\alpha_t)\,\nabla_x\log p_t(\sin(\alpha_t) y)
  \;=\; +\sin(\alpha_t)\,\theta_t^\top \nabla\varphi\big(\sin(\alpha_t)\,y\big).
  \label{eq:score-rescaled}
\end{equation}
The extra $\sin(\alpha_t)$ factor is the chain-rule term the derivation asked to
double-check -- confirmed, with the $(+)$ sign from Eq.~\ref{eq:score-sign} (a
paper-convention derivation would carry a minus here).

\emph{Practical note.} $\sin(\alpha_t)\,y = \sin(\alpha_t)\,x + \cos(\alpha_t)\,\xi$
for $y = x+\sqrt\tau\,\xi$ (since $\sqrt\tau = \cot\alpha_t$) -- literally
\code{compute\_interpolant}'s own Cos schedule with $x$ as a single-point data
endpoint and $\xi$ as fresh noise. So every $\varphi/\nabla\varphi$
evaluation below stays inside the numerical range the fit itself already
explores; nothing is evaluated off-distribution in an uncontrolled way.

\section{Denoiser form}

Substituting Eq.~\ref{eq:score-rescaled} into Tweedie's formula,
$x - \E[x\mid y] = -\sqrt\tau\,\xi - \tau\,\nabla_y\log q_\tau(y)$, so the
$1/\tau$ in Eq.~\ref{eq:denoiser-def} cancels analytically:
\begin{equation}
  \boxed{\ \deff(x,\tau) \;=\; \E_\xi\Big[\ \big\lVert\, \xi + \sqrt\tau\,\sin(\alpha_t)\,
  \theta_t^\top \nabla\varphi\big(\sin(\alpha_t)\,y\big) \,\big\rVert^2\ \Big], \qquad
  y = x + \sqrt\tau\,\xi. \ }
  \label{eq:deff-denoiser}
\end{equation}
No $\log Z_t$, no $m_t$, no normalisation constant -- only $\theta_t$ and
$\nabla\varphi$. This is the primary estimator
(\code{effective\_dimension.denoiser\_d\_eff}).

\section{Energy form}
\label{sec:energy-bound}

From $U(y,\tau) = -\log q_\tau(y)$,
\begin{equation}
  U(y,\tau) \;=\; -\theta_t^\top\varphi\big(\sin(\alpha_t) y\big) \;+\; \log Z_t
  \;-\; d\log\sin(\alpha_t).
  \label{eq:energy}
\end{equation}
The last term is analytic in $t$; the first two are not, so
Eq.~\ref{eq:energy-def}'s $\partial_s$ is taken numerically, by
finite-differencing $\E_y[U(y,\tau)\mid x]$ directly across the
$t_{\mathrm{reg}}$ grid (converted to $s=\tfrac12\log\tau(t)$), rather than
splitting analytic/Monte-Carlo pieces
(\code{effective\_dimension.energy\_d\_eff}).

\paragraph{$\log Z_t$: a bound, not an exact value.} MGD's Prop.~4.3/Eq.~30
gives, in this code's convention,
\begin{equation}
  \log Z_t \;\gtrsim\; H_{\mathrm{bound}}(t) + \theta_t^\top m_t, \qquad
  H_{\mathrm{bound}}(t) = H(p_0) + \int_0^t dH_s\,ds,
  \label{eq:logZ-bound}
\end{equation}
generalised from MGD's stated $t=1$ endpoint to any $t$ (the integral is
just truncated earlier). This is a \emph{lower} bound on $\log Z_t$, so
\textbf{the energy-form $\deff$ inherits that bound's bias} -- it is a
secondary cross-check on the denoiser form, not an independent estimate.
$H_{\mathrm{bound}}(t)$ is built by trapezoid-integrating the already-saved
$dH_t^{\mathrm{bound}}$ (\code{effective\_dimension.h\_bound\_at}); $m_t$ is
either the saved \code{aux\_moments} (the run's own realised target moments)
or, if not pulled from the cluster yet, recomputed directly from the data
and schedule, $m_t = \E_{Z}[\varphi(\cos\alpha_t\,Z + \sin\alpha_t\,x_1)]$ --
independent of any saved particle trajectory, so no SDE rerun either way
(\code{effective\_dimension.recompute\_m\_t}).

\paragraph{Empirical confirmation (2026-09-24, real 50-seed $d{=}256$
sweep).} The bound-bias prediction above was checked against a real,
full-resolution run (not a synthetic test): 50 seeds, real \code{aux\_moments}
pulled from the cluster, the full $t_{\mathrm{reg}}$ grid ($197$ nodes), 128
Monte-Carlo draws per node. Two signatures both hold, matching what a
one-sided lower bound predicts (rather than the sign-inconsistent, chaotic
pattern a residual implementation bug would produce):
\begin{itemize}
  \item \textbf{Consistent sign.} The energy form comes out \emph{lower}
        than the denoiser form in $74\%$ of all (seed, point, $\tau$)
        comparisons overall, reaching $100\%$ at the coarsest end of the
        grid.
  \item \textbf{Growing magnitude.} The median gap grows from about $-15$
        near $\tau\approx 50$ to about $-87$ near $\tau\approx 1.2\times
        10^{-3}$, with its spread growing alongside (median
        $\mathrm{std}$ across seeds rising from $\sim 5$ to $\sim 100$+ over
        the same range) -- consistent with the bound's slack accumulating
        as an (unlogged) Fisher-divergence-type residual along the
        transport, which generically grows with $t$.
\end{itemize}
Since the residual reflects how well the fitted score $\theta_t^\top
\nabla\varphi$ matches the true score at each step, \textbf{a better-converged
fit (finer $t$-grid, more potentials, better regularisation, \ldots) should
directly tighten $H_{\mathrm{bound}}(t)$ and shrink this gap} -- it is a
property of the fit's quality, not a fixed offset.

\section{Validation against a closed-form Gaussian test}
\label{sec:gaussian-test}

To separate genuine implementation bugs from the expected bound bias above,
both estimators were checked against an exact, independently-derived ground
truth -- no SDE, no fitting, no $\theta_t$ from any real run. For
$x\sim\mathcal N(0,\Sigma)$, $y=x+\sqrt\tau\,\xi$, $q_\tau=\mathcal
N(0,\Sigma+\tau I)$ exactly, so Tweedie's formula gives (diagonal basis,
$\Sigma=\mathrm{diag}(\sigma_i^2)$):
\begin{equation}
  \deff(x,\tau) \;=\; \sum_i \frac{\tau x_i^2 + \sigma_i^4}{(\sigma_i^2+\tau)^2},
  \label{eq:gaussian-oracle}
\end{equation}
verified against brute-force $(\Sigma+\tau I)^{-1}$ linear algebra (agrees
to $\sim 10^{-16}$) and its $\tau\to0$ ($\to d$) and isotropic-expectation
($\to d\sigma^2/(\sigma^2+\tau)$) limits. For \emph{this} synthetic data MGD's
own $p_t$ is also exactly Gaussian, so $\theta_t=-d/(2v(t))$ and $\log
Z_t=\tfrac{d}{2}\log(2\pi v(t))$, $v(t)=\cos^2\alpha_t+\sin^2\alpha_t\,
\sigma_{\mathrm{data}}^2$, are exact too (single-coefficient potential
$\varphi(x)=\overline{x^2}$) -- feeding these directly into the
\emph{unmodified} \code{denoiser\_d\_eff}/\code{energy\_d\_eff} isolates
implementation bugs with zero fitting noise and zero bound bias.

Using the real production grid ($d{=}256$, $t_{\mathrm{reg}}$ reconstructed
from \code{nt=40000}, \code{n\_subsample=200}): \textbf{denoiser\_d\_eff
matched the oracle to $<0.5\%$ relative error everywhere} ($10$ decades of
$\tau$, $|\theta|$ up to $512$) -- its implementation is correct.
\code{energy\_d\_eff}, despite the exact inputs, showed errors up to
$1320\%$, concentrated at the \emph{coarsest} end of the grid.

\subsection{Root cause: $s=\tfrac12\log\tau$ is highly non-uniform, and it
was differentiated naively}

$\tau(t)=\cot^2(\alpha_t)$ diverges (logarithmically, in $s$) at both
$t\to0$ and $t\to1$, so $s(t)$ is intrinsically badly stretched near the
ends of the range \emph{regardless of the $t$-schedule} -- confirmed
directly: a perfectly \emph{uniform}-in-$t$ grid has an even worse
max/min local-$\Delta s$ ratio ($340\times$) than the actual Cos-schedule
production grid ($67\times$, since bunching points near $t=1$ partially
compensates). Differentiating $\E_y[U]$ across this grid with a raw
finite difference (\code{numpy.gradient}, which defaults to
\code{edge\_order=1}) has two compounding failure modes, both confirmed in
isolation:
\begin{itemize}
  \item[(A)] \textbf{Boundary truncation bias.} \code{edge\_order=1} is
        only first-order accurate at the array's first/last nodes -- on a
        noiseless synthetic quadratic in $s$, this alone produces a
        $6.4\%$ error at the boundary (exact at every interior node).
  \item[(B)] \textbf{Monte-Carlo noise amplification} (the dominant
        effect): $U(t)$ is itself MC-estimated, and dividing its noise by
        a small local $\Delta s$ amplifies it. An $n_{\mathrm{mc}}$ sweep
        at the worst node showed the error shrinking monotonically as
        $n_{\mathrm{mc}}$ grew ($0.77\to0.29\to0.23$ over
        $4$k$\to$16k$\to$32k), confirming noise, not just truncation.
\end{itemize}
The natural first fix, \code{edge\_order=2}, removes (A) exactly (boundary
error drops to machine precision on the noiseless check) but \textbf{does
not help, and can make the realistic case worse} (max relative error
$17.9$ vs.\ $10.4$) -- its wider stencil has larger coefficients under
uneven spacing, amplifying (B) more than it fixes (A).

\subsection{Fix: a local polynomial derivative}

\code{energy\_d\_eff} now differentiates $\E_y[U]$ via
\code{effective\_dimension.local\_poly\_derivative}: at each node, a
weighted-least-squares quadratic is fit through a window of $2h+1$
\emph{neighbouring indices} (shifted, not shrunk, at the array edges, so
every node -- including the boundary -- gets a full window) and
differentiated analytically. This fixes (A) exactly (machine precision on
the noiseless check, any window size) and reduces (B) by averaging over
several points rather than two. Validated on the Gaussian test
($h=5$, the shipped default): overall median relative error dropped
$\sim6\times$ ($1.4\%\to0.22\%$), and the fraction of the grid accurate to
$<1\%$ roughly doubled ($39\%\to70\%$ of nodes, with the ``becomes
reliable'' point moving from $t\approx0.80$ down to $t\approx0.41$). The
few nodes at the very coarsest end of the grid remain the least reliable
part of an energy-form curve on principle -- no local model, however
fit, is well-conditioned there -- and should be read with that in mind;
the denoiser form has no such issue, since it involves no differentiation
step at all.

\section{Dictionary}

\begin{center}
\begin{tabular}{@{}ll@{}}
\toprule
Guth et al. & MGD (this codebase) \\
\midrule
VE noise scale $\tau$ & $\tau(t) = \cot^2(\alpha_t)$, $\alpha_t=\pi t/2$ \\
score $\nabla_y\log q_\tau(y)$ & $+\sin(\alpha_t)\,\theta_t^\top\nabla\varphi(\sin(\alpha_t) y)$ \\
energy $U(y,\tau)$ & $-\theta_t^\top\varphi(\sin(\alpha_t) y) + \log Z_t - d\log\sin(\alpha_t)$ \\
$\log Z_\tau$ & $H_{\mathrm{bound}}(t) + \theta_t^\top m_t$ \emph{(bound, Eq.~\ref{eq:logZ-bound})} \\
noised sample $y=x+\sqrt\tau\,\xi$ & $\sin(\alpha_t)\,y = $ \code{compute\_interpolant}$(x,\xi)$ \\
per-$\tau$ fit & $\theta_t \equiv$ \code{Theta\_reg[i]}, grid $t_{\mathrm{reg}}$ \\
\bottomrule
\end{tabular}
\end{center}

\section{What's used from a saved run}

\begin{itemize}
  \item \code{Theta\_reg}, shape $(n_t, r)$ -- the regularised $\theta_t$ trajectory
        (\code{codes/regularised\_theta.py}); used exclusively (not the raw,
        finer $\theta_t$) for a single self-consistent grid and to avoid a
        pre-fix off-by-one alignment issue present in some saved runs' raw
        $\theta_t$/\code{dH\_t\_bound} pairing.
  \item $t_{\mathrm{reg}}$ -- \code{Theta\_reg}'s own coarse time grid; not saved
        for the $d{=}256$ runs on disk as of this writing, so it is
        reconstructed offline (\code{reconstruct\_t\_reg}) by replaying
        \code{SDE.forward\_regularised}'s block-accumulation and
        near-node-pruning exactly, from the saved fine grid $t$ alone --
        verified to reproduce the on-disk \code{Theta\_reg} length exactly.
  \item Each run's \code{fitted\_potentials/} -- $\varphi,\nabla\varphi$ with the
        SAME pruned active-statistic set the fit used (some potentials
        auto-prune during \code{fit()}; skipping this step silently
        mis-indexes every $\theta$ coefficient).
  \item \code{dH\_t\_bound} + the fine grid $t$ -- integrated to $H_{\mathrm{bound}}(t)$
        (Eq.~\ref{eq:logZ-bound}), needed only by the energy form.
  \item \code{aux\_moments} (\code{barphi\_e}) if pulled from the cluster, else
        $m_t$ recomputed from the run's own $x_1$ -- needed only by the
        energy form.
\end{itemize}

\section{Caveats}

\begin{itemize}
  \item \textbf{The energy form's remaining disagreement with the denoiser
        form is, at this point, believed to be the expected bound bias}
        (Section~\ref{sec:energy-bound}'s empirical confirmation), not a
        residual implementation bug -- Section~\ref{sec:gaussian-test}'s
        controlled test found and fixed a genuine differentiation bug, and
        what is left afterward is sign-consistent and grows smoothly with
        $t$, the signature the bound predicts rather than of noise/bugs.
        This conclusion rests on one 50-seed sweep at one $(\sigma,\lambda)$
        setting, not a systematic sweep -- treat it as the working
        hypothesis, not a closed case.
  \item $\theta_t$ becomes numerically extreme as $t\to1$ ($\tau\to0$): every
        seed inspected so far has one dominant, near-null-direction
        coefficient at the finest scale. This is harmless for the actual
        energy/score (both stay smooth off-support -- checked directly),
        but $\theta_t^\top\varphi$ and $\theta_t^\top\nabla\varphi$ must be
        computed in double precision, not the float32 the values are stored
        in. The Gaussian test only exercised $|\theta|$ up to $512$, well
        below the real runs' $\sim10^{16}$ near $t\to1$ -- this specific
        extreme-magnitude regime is \emph{not} covered by the validation in
        Section~\ref{sec:gaussian-test} and remains a plausible contributor
        to disagreement at the very smallest $\tau$, not yet isolated from
        the bound bias there.
  \item The finest available $\tau$ is set by how close to $t=1$ the run's
        schedule and \code{n\_subsample} block reach, not by an arbitrary
        choice -- see \code{codes/effective\_dimension.py}'s
        \code{reconstruct\_t\_reg} for the exact grid actually available for
        a given run.
\end{itemize}

\section*{Implementation changelog}

\begin{itemize}
  \item \emph{2026-09-23:} \code{denoiser\_d\_eff}/\code{energy\_d\_eff} cast
        $\theta_t$ to \code{x.dtype} for \code{potentials.grad()} but not to
        \code{x.device} -- harmless on a single-device (CPU) run, but a hard
        crash on Jean Zay, where $\theta_t$/$m_t$ load onto CPU
        (\code{map\_location='cpu'}) while \code{x}/the potentials live on
        \code{cuda}. Fixed: both are now moved to \code{x.device} inside
        each function.
  \item \emph{2026-09-24:} the differentiation bug and fix described in
        Section~\ref{sec:gaussian-test} (\code{local\_poly\_derivative}
        replacing a raw \code{numpy.gradient} call in
        \code{energy\_d\_eff}).
\end{itemize}

\end{document}
